% Generated draft: requires mathematical and source review.
\documentclass[11pt]{article}
\usepackage[margin=25mm]{geometry}
\usepackage{fontspec}
\setmainfont{Latin Modern Roman}
\usepackage{amsmath,amssymb,mathtools,unicode-math}
\setmathfont{Latin Modern Math}
\usepackage{enumitem,needspace,xurl,hyperref}
\setlength{\parindent}{0pt}
\setlength{\parskip}{5pt}
\setlength{\emergencystretch}{3em}
\newcommand{\N}{\mathbb{N}}
\newcommand{\Z}{\mathbb{Z}}
\newcommand{\Q}{\mathbb{Q}}
\newcommand{\R}{\mathbb{R}}
\newcommand{\C}{\mathbb{C}}
\newcommand{\F}{\mathbb{F}}
\newcommand{\A}{\mathbb{A}}
\newcommand{\PP}{\mathbb{P}}
\newcommand{\E}{\mathbb{E}}
\newcommand{\eps}{\varepsilon}
\newcommand{\dd}{\,\mathrm{d}}
\begin{document}
\section*{UnsolvedMath ID: 20000944}
\textbf{Source problem code:} catalogue entry
\subsection{Definitions and mathematical statement}
A graph is $d$-degenerate if every nonempty induced subgraph has a vertex of degree at most $d$. For a $k$-uniform hypergraph $H$, ordinary hypergraph degeneracy uses incident-hyperedge degree. Its clique shadow $\partial_2H$ is the graph joining two vertices whenever they occur together in a hyperedge. Call $H$ skeletally $s$-degenerate if $\partial_2H$ is $s$-degenerate. These two hypergraph notions are distinct.

Let $r_k(H,H)$ be the least $N$ such that every red/blue coloring of the complete $k$-graph on $N$ vertices contains a monochromatic copy of $H$.

\textbf{Source item 22.} The Burr--Erd\H{o}s conjecture asks whether every $n$-vertex $d$-degenerate graph satisfies $r_2(G,G)\leq c(d)n$. The source records the earlier Kostochka--Sudakov bound $n^{1+o_d(1)}$. It says the analogous assertion can fail for ordinary-degenerate $3$-graphs and gives an ordinary-$1$-degenerate $4$-uniform family with
\[
 r_4(H,H)\geq2^{\Omega(n^{1/3})}.
\]
It then proposes replacing hypergraph degeneracy by degeneracy of the clique shadow, namely the skeletal definition above.

\textbf{Supplied resolution.} The aid states that Lee's 2017 theorem proves the graph conjecture. It also credits Fox--Sankar--Simkin--Tidor--Zhou, Theorem 1.5, with the hypergraph linear-Ramsey analogue for exactly the proposed skeletal notion: with $k,s$ fixed, $r_k(H,H)\leq C(k,s)|V(H)|$ for skeletally $s$-degenerate $k$-graphs. This does not assert a linear bound for all ordinary-degenerate hypergraphs. Complete citations and precise theorem hypotheses are not supplied.

\textbf{Reported sharp structural consequences.} Let $H$ be a simple $k$-uniform hypergraph, $k\geq2$, whose skeletal degeneracy is at most $s$. If it has an edge then $s\geq k-1$. For every vertex subset $U$ with $m=|U|$ the aid proves
\[
 |E(H[U])|\leq
 \begin{cases}
 \binom mk,&m<s,\\
 \binom sk+(m-s)\binom{s}{k-1},&m\geq s.
 \end{cases}
\]
It also proves ordinary hypergraph degeneracy at most $\binom{s}{k-1}$. The shadow of $H[U]$ is a subgraph of the restricted shadow, so the skeletal bound is inherited even though those shadows need not be equal.

The reported edge envelope is sharp in the class with skeletal degeneracy at most $s$: take all $k$-cliques of the graph formed by joining $K_s$ to an independent set. For $m<s$, the complete $k$-graph attains the small-set bound. The complete $k$-graph on $s+1$ vertices has skeletal degeneracy $s$ and ordinary degeneracy $\binom{s}{k-1}$, attaining the latter bound. The structural statements relate the two sparsity notions but do not replace the cited Ramsey proof.
\subsection{Short English statement}
The source proposes clique-shadow degeneracy to repair failed hypergraph analogues of Burr--Erd\H{o}s. The aid reports that the graph conjecture and this skeletal hypergraph version are proved in the literature. It adds sharp bounds on every induced hypergraph's edge count and on its ordinary degeneracy, keeping the two notions separate.
\subsection{Sources}
\begin{itemize}
\item Local UnsolvedMath registry, permanent ID $20000944$, \texttt{catalogue entry}; canonical \texttt{aim-combinatorics-notes.json}, \texttt{notes[68]}.
\item AIM workshop ``Graph Ramsey theory,'' section ``Chromatic number of graphs with everywhere high independence number,'' item 22: \url{https://www.overleaf.com/read/mnvcscjjysvg}.
\item Original names Burr--Erd\H{o}s and Kostochka--Sudakov. The research summary names Lee's 2017 theorem and Fox--Sankar--Simkin--Tidor--Zhou, Theorem 1.5, but gives no full titles or identifiers.
\item Research attempt 1, \texttt{partial\_result}, run status \texttt{solved}, internal rigor \texttt{proved}, low novelty confidence; registry \texttt{SOLVED-IN-LITERATURE}. Explicitly machine-generated, not peer reviewed, and to be independently verified before citation. Original literature field empty.
\end{itemize}
\end{document}
